\section{Appendix to Section \ref{app:eff}}\label{app:app:eff}

\begin{proof}[Additional details for the proof of Lemma \ref{lemma:pca-pcm}]
    Assume that $X$ is the free pca generated by a set $S$ and recall that free pcas are cancellative. We first prove that for the structure $(X,\circledvee,\zeropcm)$ defined above $\circledvee$ is well-defined, then we prove that the axioms of pcm are satisfied.
    
    Hence, assume that $x\bot y$, which means that there exist $x',y'\in X$ and $p,q\in [0,1]$ such that $x=p\cdot x'$, $y=q\cdot y'$ and $p+q\leq 1$. Assume also that $x\bot y$ via $x'',y''\in X$ and $r,s\in [0,1]$ such that $x=r\cdot x''$, $y=s\cdot y''$ and $r+s\leq 1$. First consider the case where $p,q,r,s\in(0,1)$, then we need to prove the equality
\begin{equation}\label{eq:pcmbendefinita}
    x'+_p(y'+_{\frac{q}{1-p}}\star) = x''+_r(y''+_{\frac{s}{1-r}}\star)\text{.}\end{equation}
In order to do that, we assume that $p\leq r$ and that $s\leq q$ (the other cases can be treated similarly).
First observe that 
\begin{align*}
    (1-p)\cdot (\frac{q}{1-p}\cdot y') & = q\cdot y' \\
    &= s\cdot y'' \tag{Hypothesis}\\
    & = (1-p)\cdot (\frac{s}{1-p}\cdot y'')\text{.}
    \end{align*}
    Thus, by cancellativity, we obtain that 
    \begin{equation}\tag{$\dagger$}\label{pcmtemp1}
    \frac{q}{1-p}\cdot y' = \frac{s}{1-p}\cdot y''.
    \end{equation}
     One can similarly show that 
     \begin{equation}\tag{$\ddagger$}\label{pcmtemp2}
     \frac{p}{1-s}\cdot x'=\frac{r}{1-s}\cdot x''.
     \end{equation}
    Now we show that both sides of \eqref{eq:pcmbendefinita} are equal to $x'+_p(y''+_{\frac{s}{1-p}}\star)$. We have %\marginpar{Perche' si riferisce a 2.2???}
\begin{align*}
    x'+_p(y'+_{\frac{q}{1-p}}\star) & = x'+_p((\frac{q}{1-p}\cdot y')) \tag{Def.\ of $p\cdot x$}\\
    & = x'+_p((\frac{s}{1-p}\cdot y'')) \tag{\ref{pcmtemp1}}\\
    & = x'+_p(y''+_{\frac{s}{1-p}}\star) \tag{Def.\ of $p\cdot x$}
    \end{align*}
    Similarly,
    \begin{align*}
 x''+_r(y''+_{\frac{s}{1-r}}\star) & = y''+_s (x''+_{\frac{r}{1-s}}\star) \tag{using axioms in \ref{eq:pca}}\\
    & = y''+_s(\frac{r}{1-s}\cdot x'') \tag{Def.\ of $p\cdot x$}\\
    & = y''+_s(\frac{p}{1-s}\cdot x') \tag{\ref{pcmtemp2}}\\
    & = y''+_s(x'+_{\frac{p}{1-s}}\star) \tag{Def.\ of $p\cdot x$}\\
    & = x'+_p(y''+_{\frac{s}{1-p}}\star ). \tag{using axioms in \ref{eq:pca}}
\end{align*}
Now assume that $p=1$ and $q=0$, then $x\circledvee y = x$. We distinguish two subcases: either $r=1$ or $r\not=1$. In the first case, the claim holds trivially. In the second case, we have that $q=0$ and hence, if $s\not= 0$, we have $s\cdot y''= q\cdot y'= 0 \cdot y' = \star = s\cdot \star$ and then, by cancellativity, $y''=\star$. Thus, $x''+_r(y''+_{\frac{s}{1-r}}\star) = x''+_r \star = r\cdot x'' = x$. The remaining cases can be treated similarly. Hence the equality in \eqref{eq:pcmbendefinita} holds and $\circledvee$ is well defined.


Now we prove the axioms of pcm. To prove commutativity, assume that $x\bot y$, and hence there exist $x',y'\in X$ and $p,q\in [0,1]$ such that $x=p\cdot x'$, $y=q\cdot y'$, $p+q\leq 1$. Then $y\bot x$ via $y',x'$ and $q,p$. 
Moreover, if $p,q\not=1$, then
\begin{align*}
    x\circledvee y & = x'+_p(y'+_{\frac{q}{1-p}}\star) \tag{\eqref{eq:defcircledvee}}\\
    & = (y'+_{\frac{q}{1-p}}\star) +_{1-p} x' \tag{symmetry in \eqref{eq:pca}}\\
    & = y'+_q(\star +_{\frac{1-p-q}{1-q}} x') \tag{associativity in \eqref{eq:pca}}\\
    & = y'+_q(x'+_{\frac{p}{1-q}}\star) \tag{symmetry in \eqref{eq:pca}}\\
& = y\circledvee x. \tag{\eqref{eq:defcircledvee}}
\end{align*}
If $p=1$ and $q=0$, then $x\circledvee y = x = y\circledvee x$. The case $p=0$ and $q=1$ is similar.


To prove that $\zeropcm \bot x$ and $\zeropcm \circledvee x =x$, we observe that $\zeropcm=\star = 0\cdot \star$ and $x=1\cdot x$, and hence by \eqref{eq:defbot}, $\zeropcm \bot x$. Obviously, $\zeropcm \circledvee x = x +_1 \star = x$. 

Associativity has been already proved in the main text.
\end{proof}